% SPDX-License-Identifier: CC-BY-SA-4.0
% Part: Lexical Analysis
% Section: Deterministic and complete automata
% Exercise: Determinisation using Rabin-Scott's algorithm

\begingroup

\begin{exercice}[Propriétés des automates]\label{exo:lexing/determinisation/determinisation}

  Pour chacun des automates suivant définis sur l'alphabet $\{a,b\}$ : 
  \begin{itemize}
  \item Dire s'il est unitaire, standard, normalisé, $\varepsilon$-libre, déterministe et complet.
  \item Si l'automate n'est pas complet, le compléter. 
  \item Si l'automate n'est pas déterministe, donner l'automate déterministe équivalent par la méthode de construction des 
    sous-ensembles de Rabin et Scott.
  \end{itemize}

  \begin{question}

  \item 
    \begin{tikzpicture}[automaton, size=15mm, baseline=(0.base)]
      \state[initial]   (0) at (0,0) {$0$}; 
      \state            (1) at (1,0) {$1$}; 
      \state[accepting] (2) at (2,0) {$2$}; 
      \path (0) edge             node {$a$} (1);
      \path (1) edge             node {$b$} (2);
      \path (2) edge[loop above] node {$a$} (2);
    \end{tikzpicture}    

  \item \begin{tikzpicture}[automaton, size=15mm, baseline=(0.base)]
    \state[initial]   (0) at (0,0) {$0$}; 
    \state            (1) at (1,0) {$1$}; 
    \state[accepting] (2) at (2,0) {$2$}; 
    \path (0) edge[loop above] node {$a, b$}  (0);
    \path (1) edge[loop above] node {$a, b$}  (1);
    \path (2) edge[loop above] node {$a, b$}  (2);
    \path (0) edge             node {$a$}     (1);
    \path (1) edge             node {$a$}     (2);
  \end{tikzpicture}
    
  \item \begin{tikzpicture}[automaton, size=15mm, baseline=(1.base)]
    \state[initial]   (1) at (0,1) {$1$}; 
    \state            (2) at (1,2) {$2$}; 
    \state            (3) at (1,0) {$3$}; 
    \state[accepting] (4) at (2,1) {$4$}; 

    \path (1) edge[bend left] node {$a$} (2);
    \path (1) edge            node {$b$} (3);
    \path (2) edge            node {$a$} (1);
    \path (2) edge[bend left] node {$b$} (4);
    \path (3) edge[bend left] node {$b$} (1);
    \path (3) edge            node {$a$} (4);
    \path (4) edge            node {$b$} (2);
    \path (4) edge[bend left] node {$a$} (3);
  \end{tikzpicture}

          \item 
    \begin{tikzpicture}[automaton, size=15mm, baseline=(1.base)]
      \state[initial, accepting] (1) at (0,1) {$1$}; 
      \state                     (2) at (1,1) {$2$}; 
      \state                     (3) at (2,1) {$3$}; 
      \state                     (4) at (0,0) {$4$}; 
      \state                     (5) at (1,0) {$5$}; 
      \state[accepting]          (6) at (2,0) {$6$}; 

      \path (1) edge[loop above] node {$a$}   (1);
      \path (4) edge[loop below] node {$a,b$} (4);
      \path (5) edge[loop below] node {$b$}   (5);
      \path (1) edge             node {$b$}   (2);
      \path (1) edge             node {$a$}   (4);
      \path (2) edge             node {$a$}   (3);
      \path (2) edge             node {$b$}   (4);
      \path (3) edge[bend left]  node {$a$}   (6);
      \path (6) edge[bend left]  node {$a$}   (3);
      \path (6) edge             node {$a$}   (5);
    \end{tikzpicture}

  \item 
    \begin{tikzpicture}[automaton, size=15mm, baseline=(1.base)]
      \state[initial]   (1)  at (0,2) {$1$}; 
      \state            (2)  at (1,2) {$2$}; 
      \state            (3)  at (0,1) {$3$}; 
      \state            (4)  at (2,2) {$4$}; 
      \state            (5)  at (3,2) {$5$}; 
      \state            (6)  at (4,2) {$6$}; 
      \state            (7)  at (1,1) {$7$}; 
      \state            (8)  at (2,1) {$8$}; 
      \state[initial]   (9)  at (0,0) {$9$}; 
      \state            (10) at (1,0) {$10$}; 
      \state[accepting] (11) at (4,1) {$11$}; 

      \path (1)  edge            node[swap] {$\varepsilon$} (2);
      \path (2)  edge            node[swap] {$\varepsilon$} (4);
      \path (4)  edge            node       {$a$}           (5);
      \path (5)  edge            node[swap] {$\varepsilon$} (6);
      \path (6)  edge            node       {$\varepsilon$} (11);
      \path (1)  edge            node       {$\varepsilon$} (3);
      \path (3)  edge            node       {$a$}           (7);
      \path (7)  edge            node       {$b$}           (8);
      \path (8)  edge            node       {$\varepsilon$} (11);
      \path (9)  edge            node       {$\varepsilon$} (10);
      \path (10) edge            node[swap] {$b$}           (8);
      \path (2)  edge[bend left] node       {$\varepsilon$} (6);
      \path (5)  edge[bend left] node       {$\varepsilon$} (4);
    \end{tikzpicture}
    
  \end{question}

\end{exercice}

\begin{correction}

  \textbf{Rappels de vocabulaire.}
  \begin{description}
  \item[Unitaire] : un seul état initial.
  \item[Standard] : unitaire sans transition entrante dans l'état initial.
  \item[Normalisé] : standard + un seul état final sans transition sortante.
  \item[$\varepsilon$-libre] : aucune $\varepsilon$-transition.
  \item[Déterministe] : Unitaire, $\varepsilon$-libre, et au plus une transition par lettre et par état.
  \item[Complet] : au moins une transition par lettre et par état.
  \end{description}
  
  \begin{question}

  \item
    \begin{itemize}
    \item L'automate est : unitaire, standard, $\varepsilon$-libre, déterministe.
    \item L'automate n'est pas : normalisé, complet.
    \item Complétion :
      \begin{tikzpicture}[automaton, size=15mm, baseline=(0.base)]
        \state[initial]   (0) at (0,1) {$0$}; 
        \state            (1) at (1,1) {$1$}; 
        \state[accepting] (2) at (2,1) {$2$}; 
        \state            (3) at (1,0) {$\bot$}; 
        \path (0) edge             node {$a$}   (1);
        \path (1) edge             node {$b$}   (2);
        \path (2) edge[loop above] node {$a$}   (2);
        \path (0) edge             node {$b$}   (3);
        \path (1) edge             node {$a$}   (3);
        \path (2) edge             node {$b$}   (3);
        \path (3) edge[loop below] node {$a,b$} (3);
      \end{tikzpicture}
    \end{itemize}

  \item 
    \begin{itemize}
    \item L'automate est : unitaire, $\varepsilon$-libre, complet.
    \item L'automate n'est pas : standard, normalisé, déterministe.
    \item Déterminisation : On donne uniquement les détails de l'algorithme pour cette question.
      
      \emph{Phase 1.}
      La construction de l'AFD équivalent commence par la détermination de l'état initial.
      Celui-ci est construit à partir de l' $\varepsilon$-fermeture de tous les états initiaux.
      Donc ici, $I'=\{0\}$. A partir de là, on construit la table des transitions.
      
      $$\begin{array}{| l | c | c|}
        \hline
        \  \mu & a & b \\
        \hline
        \  A = \{0\} = I' &  &  \\
        \hline
      \end{array}$$
      
      \emph{Phase 2.}
      On regarde vers quels états l'automate se trouve pour chacun des symboles de l'alphabet à partir de $I'$.
      Normalement, il faut faire l'$\varepsilon$-fermeture sur les états obtenus mais dans cet automate il n'y a aucune $\varepsilon$-transition.
      On les ajoute à la table s'ils n'y sont pas déjà.
      
      $$\begin{array}{| l | c | c|}
        \hline
        \  \mu & a & b \\
        \hline
        \  A = \{0\} = I' & \{0,1\} & \{0\} \\
        \  B = \{0,1\} &  &  \\
        \hline
      \end{array}$$
      
      
      \emph{Phase 3.}
      Puis on fait la même chose pour chacun des états pas encore traités.
      $$\begin{array}{| l | c | c|}
        \hline
        \  \mu & a & b \\
        \hline
        \  A = \{0\} = I' & \{0,1\} & \{0\} \\
        \  B = \{0,1\} &  \{0,1,2 \}& \{0,1\}  \\
        \  C = \{0,1,2\} &  &  \\
        \hline
      \end{array}$$
      
      \emph{Phase 4.}
      C contient l'état 2 qui est final dans l'automate de départ. Donc C est un état final dans l'AFD obtenu.
      
      $$\begin{array}{| l | c | c|}
        \hline
        \  \mu & a & b \\
        \hline
        \  A = \{0\} = I' & \{0,1\} & \{0\} \\
        \  B = \{0,1\} &  \{0,1,2 \} & \{0,1\}  \\
        \  C = \{0,1,2\} \in F &  \{0,1,2 \} &  \{0,1,2 \} \\
        \hline
      \end{array}$$
      
      Finalement, le tableau est rempli, l'automate est terminé. On obtient l'automate suivant :
      
      \begin{center}
        \begin{tikzpicture}[automaton]
          \state[initial  ] (0) at (0,0) {$A$}; 
          \state            (1) at (1,0) {$B$}; 
          \state[accepting] (2) at (2,0) {$C$}; 
          
          \path (0) edge[loop above] node {$b$}    (0);
          \path (1) edge[loop above] node {$b$}    (1);
          \path (2) edge[loop above] node {$a, b$} (2);
          \path (0) edge             node {$a$}    (1);
          \path (1) edge             node {$a$}    (2);
        \end{tikzpicture}
      \end{center}
    \end{itemize}
    
  \item
    \begin{itemize}
    \item L'automate est : unitaire, $\varepsilon$-libre, déterministe, complet.
    \item L'automate n'est pas : standard, normalisé.
    \end{itemize}

  \item 
    \begin{itemize}
    \item L'automate est : unitaire, $\varepsilon$-libre.
    \item L'automate n'est pas : standard, normalisé, déterministe, complet.
    \item Automate déterministe et complet :

      \begin{tikzpicture}[automaton, baseline=(a.base), x=20mm]
        \state[accepting]          (a) at (0,3) {$14$};
        \state                     (b) at (1,3) {$24$};
        \state                     (c) at (2,3) {$34$};
        \state[accepting]          (d) at (3,3) {$64$};
        \state                     (e) at (4,3) {$345$};
        \state                     (f) at (4,2) {$45$};
        \state                     (g) at (2,2) {$4$};
        \state[initial, accepting] (h) at (0,1) {$1$};
        \state                     (i) at (1,1) {$2$};
        \state                     (j) at (2,1) {$3$};
        \state[accepting]          (k) at (3,1) {$6$};
        \state                     (l) at (4,1) {$35$};
        \state                     (m) at (4,0) {$5$};
        \state                     (n) at (2,0) {$\bot$};
        
        \path (a) edge[loop left]  node {$a$}    (a);
        \path (a) edge             node {$b$}    (b);
        \path (b) edge             node {$a$}    (c);
        \path (b) edge             node {$b$}    (g);
        \path (c) edge             node {$a$}    (d);
        \path (c) edge             node {$b$}    (g);
        \path (d) edge[bend left]  node {$a$}    (e);
        \path (d) edge             node {$b$}    (g);
        \path (e) edge[bend left]  node {$a$}    (d);
        \path (e) edge             node {$b$}    (f);
        \path (f) edge             node {$a$}    (g);
        \path (f) edge[loop right] node {$b$}    (f);
        \path (g) edge[loop left]  node {$a, b$} (g);
        \path (h) edge             node {$a$}    (a);
        \path (h) edge             node {$b$}    (i);
        \path (i) edge             node {$a$}    (j);
        \path (i) edge             node {$b$}    (g);
        \path (j) edge             node {$a$}    (k);
        \path (j) edge             node {$b$}    (n);
        \path (k) edge[bend left]  node {$a$}    (l);
        \path (k) edge             node {$b$}    (n);
        \path (l) edge[bend left]  node {$a$}    (k);
        \path (l) edge             node {$b$}    (m);
        \path (m) edge             node {$a$}    (n);
        \path (m) edge[loop right] node {$b$}    (m);
        \path (n) edge[loop left]  node {$a, b$} (n);

      \end{tikzpicture}
    \end{itemize}

  \item
    \begin{itemize}
    \item L'automate n'est pas : unitaire, standard, normalisé, $\varepsilon$-libre, déterministe, complet.
    \item On pose : 
      \begin{itemize}
      \item $A = \{1, 2, 3, 4, 6, 9, 10, 11\}$
      \item $B = \{8, 11\}$
      \item $C = \{4, 5, 6, 7, 11\}$
      \item $D = \{4, 5, 6, 11\}$
      \end{itemize}

      \begin{tikzpicture}[automaton]
        \state[accepting, initial] (A) at (0,1) {$A$}; 
        \state[accepting]          (B) at (1,1) {$B$}; 
        \state[accepting]          (C) at (0,0) {$C$}; 
        \state[accepting]          (D) at (1,0) {$D$}; 

        \path (D) edge[loop right] node {$a$}  (D);
        \path (A) edge             node {$b$} (B);
        \path (A) edge             node {$a$} (C);
        \path (C) edge             node {$b$} (B);
        \path (C) edge             node {$a$} (D);
      \end{tikzpicture}
    \end{itemize}

  \end{question}

\end{correction}

\begin{correction}

  \begin{question}
  \item {\small
    \begin{tabular}[t]{| l | c c c c |}
      \hline
      & A & B & C  & D \\
      \hline
      Accessible          & $\{1,2,3,4\}$ & $\{1, 2, 3, 4, 5, 7\}$  & $\{1, 2, 3, 4, 5, 6\}$  & $\{1, 2, 3, 4, 5, 6, 7, 8, 9, 10, 11\}$ \\
      Inaccessible        & $\emptyset$   & $\{6\}$                 & $\emptyset$             & $\emptyset$                             \\
      Co-accessible       & $\{1,2,3,4\}$ & $\{1, 2, 5\}$           & $\{1, 2, 3, 6\}$        & $\{1, 2, 3, 4, 5, 6, 7, 8, 9, 10, 11\}$ \\
      Stérile             & $\emptyset$   & $\{3, 4, 6, 7\}$        & $\{4, 5\}$              & $\emptyset$                             \\
      Utile               & $\{1,2,3,4\}$ & $\{1, 2, 5\}$           & $\{1, 2, 3, 6\}$        & $\{1, 2, 3, 4, 5, 6, 7, 8, 9, 10, 11\}$ \\
      Inutile             & $\emptyset$   & $\{3, 4, 6, 7\}$        & $\{4, 5\}$              & $\emptyset$                             \\
      Complet             & oui           & non                     & non                     & non                                     \\
      Emondé              & oui           & non                     & non                     & oui                                     \\
      Unitaire            & oui           & oui                     & oui                     & oui                                     \\
      Standard            & non           & oui                     & non                     & oui                                     \\
      Normalisé           & non           & non                     & non                     & non                                     \\
      $\varepsilon$-libre & oui           & oui                     & oui                     & non                                     \\
      Déterministe        & oui           & non                     & non                     & non                                     \\
      \hline
  \end{tabular}}
  \item 
    \begin{minipage}[t]{.4\textwidth}
      $B$ sur l'alphabet $\{a,b,c,d\}$ :\\

      \scalebox{.9}{\begin{tikzpicture}[shorten >=1pt,node distance=1.7cm,on grid,auto]
          \node[state, initial, initial text=, accepting] (1)              {$1$}; 
          \node[state, accepting]                         (2) [right=of 1] {$2$}; 
          \node[state, accepting]                         (5) [below=of 2] {$5$}; 

          \path[-latex] (2) edge[loop above, looseness=5] node {$a$}    (2);
          \path[-latex] (5) edge[loop below, looseness=5] node {$b, c$}    (5);
          \path[-latex] (1) edge node {$a$} (2);
          \path[-latex] (1) edge node[swap] {$b, c$} (5);
          \path[-latex] (2) edge[bend left=6mm] node {$b$} (5);
          \path[-latex] (5) edge[bend left=6mm] node {$a$} (2);
      \end{tikzpicture}}
    \end{minipage}%
    \begin{minipage}[t]{.4\textwidth}
      $C$ sur l'alphabet $\{0,1,2,3,4,5,6,7,8,9\}$ :\\
      \scalebox{.9}{\begin{tikzpicture}[shorten >=1pt,node distance=1.7cm,on grid,auto]
          \node[state, initial, initial text=, accepting] (1)              {$1$}; 
          \node[state]                                    (2) [right=of 1] {$2$}; 
          \node[state]                                    (3) [right=of 2] {$3$}; 
          \node[state, accepting]                         (6) [below=of 3] {$6$}; 

          \path[-latex] (1) edge[loop above, looseness=5] node {$0$}    (1);
          \path[-latex] (1) edge node {$1$} (2);
          \path[-latex] (2) edge node {$0$} (3);
          \path[-latex] (3) edge[bend left=6mm] node {$0$} (6);
          \path[-latex] (6) edge[bend left=6mm] node {$0$} (3);
      \end{tikzpicture}}
    \end{minipage}
  \end{question}

\end{correction}

\endgroup
\endinput

